\documentclass[11pt,a4paper]{article}
\usepackage[T1]{fontenc}
\usepackage[utf8]{inputenc}
\usepackage{lmodern,microtype}
\usepackage[margin=25mm,headheight=15pt]{geometry}
\usepackage{amsmath,amssymb,amsthm,mathtools}
\usepackage{booktabs,longtable,tabularx,array}
\usepackage{xcolor,fancyhdr,enumitem}
\usepackage{hyperref}
\definecolor{navy}{RGB}{25,48,73}
\definecolor{muted}{RGB}{75,87,100}
\hypersetup{colorlinks=true,linkcolor=navy,citecolor=navy,urlcolor=navy,pdftitle={Correctness Audit of Boolean CM Modal Models},pdfauthor={Research audit prepared for Brian}}
\pagestyle{fancy}\fancyhf{}
\fancyhead[L]{\small\textcolor{muted}{Boolean CM: correctness and model boundaries}}
\fancyhead[R]{\small\textcolor{muted}{18 September 2026}}
\fancyfoot[C]{\thepage}
\setlength{\parindent}{0pt}\setlength{\parskip}{6pt}
\setlength{\emergencystretch}{2em}
\renewcommand{\arraystretch}{1.18}
\setlength{\tabcolsep}{5pt}
\setlist{nosep,leftmargin=1.4em}
\newcolumntype{Y}{>{\raggedright\arraybackslash}X}
\newcolumntype{L}[1]{>{\raggedright\arraybackslash}p{#1}}
\newtheorem{theorem}{Theorem}[section]
\newtheorem{proposition}[theorem]{Proposition}
\newtheorem{lemma}[theorem]{Lemma}
\theoremstyle{definition}\newtheorem{definition}[theorem]{Definition}
\newtheorem{remark}[theorem]{Remark}
\newcommand{\F}{\mathbb F_2}
\newcommand{\A}{\mathcal A}
\newcommand{\B}{\mathbb B}
\newcommand{\GL}{\operatorname{GL}}
\newcommand{\rk}{\operatorname{rank}_{\oplus,\wedge}}
\newcommand{\im}{\operatorname{im}}
\newcommand{\spn}{\operatorname{span}}
\newcommand{\vecop}{\operatorname{vec}}
\newcommand{\ket}[1]{\lvert #1\rangle}
\newcommand{\bra}[1]{\langle #1\rvert}
\newcommand{\zero}{\mathbf 0}
\newcommand{\ev}[1]{\textbf{Evidence:} #1.}
\newcommand{\file}[1]{\path{#1}}
\newcommand{\modelS}{\textsf{S}}
\newcommand{\modelL}{\textsf{L}}
\newcommand{\modelC}{\textsf{C}}
\newcommand{\modelR}{\textsf{R}}
\begin{document}
\begin{titlepage}
\vspace*{14mm}
{\small\sffamily\textcolor{muted}{QUANTUM CM RESEARCH / VERIFIED BASELINE 1.0}}\par
\vspace{9mm}
{\Huge\bfseries\textcolor{navy}{Correctness Audit of\\[3mm]Boolean CM Modal Models}}\par
\vspace{7mm}
{\Large Derivations, exhaustive finite checks,\\model boundaries, and retained failures}\par
\vspace{12mm}
\begin{tabularx}{\textwidth}{@{}L{34mm}Y@{}}
Prepared for & Brian / Quantum CM project\\
Audit date & 18 September 2026\\
Core operations & Boolean AND and XOR contraction; permutations, rotation, transpose, and explicitly typed tensor products\\
Deliverables & This PDF, editable LaTeX source, independent code, raw evidence, historical source snapshots, and reproducibility instructions
\end{tabularx}
\vspace{12mm}
\hrule\vspace{5mm}
\textbf{Principal finding.} The reversible Boolean/modal superposition, interference, Bell/GHZ, teleportation, and dense-coding constructions have valid algebraic realizations. However, they do \emph{not} all belong to one unchanged operational model. Several earlier conclusions require narrower assumptions, and some do not survive the move from a shared coefficient module to independent phase registers with unrestricted Boolean-linear gates.

\textbf{Most consequential corrections.} The coefficient-to-operator embedding is not a tensor equivalence; nonzero shared-module preparations can tensor to zero; the literal universal resource includes phase entanglement; the 64-outcome protocol uses an enlarged gate set; finite-copy no-activation fails in that enlarged model; and the Bell support table has no support-faithful no-signalling probability extension.
\vfill
{\small This is a mathematical and computational audit, not a certification of a physical quantum device, a novel quantum algorithm, or every theorem that could be formulated about CMs. No signed lift, complex amplitude, normalization, or Born rule is used in the core constructions.}
\end{titlepage}
\tableofcontents
\newpage

\section{Executive findings and the corrected baseline}
\subsection{What survives}
The most defensible operational baseline is the \emph{literal independent-register modal model}: each local system has an eight-dimensional vector space over Boolean XOR/AND arithmetic, nonzero vectors are states, composite systems use the ordinary Boolean tensor product, reversible evolution is an explicitly allowed invertible Boolean-linear map, and measurement is an added possible/impossible basis-readout rule. This is an instance of finite-field modal quantum theory, with a distinguished CM rotation and optional rotation-compatible restrictions \cite{schumacher2011,james2011}.

Within that baseline, the following constructions are valid: reversible modal superposition and XOR interference; tensor nonseparability; the tested Bell and GHZ contradictions; universal exact one-copy transfer with a full-rank correlated resource and a complete 64-outcome analyzer; and 64-message dense coding using one transmitted eight-dimensional subsystem and a shared resource. They are \emph{algebraic/modal} results. They are not consequences of truth-table notation alone.

\begin{table}[htbp]
\centering\small
\begin{tabularx}{\textwidth}{@{}L{34mm}Y L{40mm}@{}}
\toprule
Claim & Correct final statement & Status\\\midrule
CM rotation and phase algebra & The order-four permutation $P$ generates a 16-element commuting operator algebra by XOR and composition. & Proved; exhaustive centralizer scan.\\
Superposition and interference & Reversible shears create multibranch vectors and later XOR cancellation. & Proved; explicit tests.\\
Entanglement analogue & Non-simple tensors exist. Literal bipartite separability is matrix rank one, not the old ring criterion. & Proved; both classifications checked.\\
Bell/GHZ & The stated modal tables have no deterministic global assignment. & Explicit contradictions; exhaustive finite checks.\\
Universal teleportation & With full literal resource rank eight, a fixed 64-outcome Boolean analyzer and invertible corrections recover every nonzero local state. & Proved; direct tensor execution.\\
Dense coding & 64 local encodings give 64 linearly independent joint states, all decoded exactly. & Proved; 64/64 checks.\\
Physical quantum theory & Neither a physical implementation nor the probabilistic structure of quantum mechanics is established. & Not claimed; a probability obstruction is proved.\\
\bottomrule
\end{tabularx}
\caption{Positive results do not erase the operation, tensor, and measurement assumptions.}
\end{table}

\subsection{Corrections that must accompany those results}
\begin{longtable}{@{}L{43mm}L{111mm}@{}}
\toprule Earlier interpretation & Audited correction\\\midrule\endhead
The CM identity squares to zero. & False. The $2\times2$ truth pattern $\delta_{\rm CM}=I_2$ remains idempotent. The \emph{different} $4\times4$ operator $I_4\oplus P^2$ squares to zero.\\
Removing a primitive cyclic product invalidates its entire algebra. & Too strong. The cyclic rule is exactly composition of the derived rotation operators. It is not pointwise AND or native $2\times2$ composition.\\
The shared and literal constructions are merely the same theory in different coordinates. & False. Dimensions, tensor products, preparations, entanglement, and operational resources differ.\\
All 256 local vectors are states. & The zero vector is only a null regression input. There are 255 nonzero local modal states.\\
The same logical Bell pair teleports the entire literal local system. & False. A phase-anchored logical Bell pair has rank two. Universal eight-dimensional transfer uses rank eight and includes phase correlations.\\
Four Bell outcomes still suffice after independent tensoring. & False for the old logical-only analyzer. Each of its four labels leaves 16 unresolved phase sectors with distinct Bob maps.\\
Everything in the literal protocol is generated by the commuting $P$ algebra alone. & False. A shear $T$ not commuting with $P$ is used; $P,T$ generate all of $\GL(4,2)$. Joint analyzers also require allowed cross-register wiring.\\
The 15 Smith classes classify the enlarged literal theory. & Only under the restricted rotation-compatible action. Under full local $\GL(8,2)$, matrix rank classifies bipartite pure-state orbits.\\
A singular $(0,2)$ resource cannot be activated with finite copies. & True only for the earlier rotation-linear, tensor-over-$\A$ architecture. In the literal model two copies can heraldedly yield a rank-eight resource.\\
Different literal/shared contextuality tables reveal an unexplained invariant hierarchy. & Their entire restricted empirical models are related by conjugating Bob's settings and relabeling outcomes. The matching totals have an explicit explanation.\\
No-cloning is a restriction on ordinary Boolean computation. & False. It is a restriction on the chosen \emph{linear modal tensor-state} dynamics. Stored coefficient strings can be copied; nonlinear Boolean coefficient maps can construct their tensor square.\\
Only a Born rule remains to be supplied. & Too optimistic. No no-signalling probability distribution can reproduce the tested Bell table with positive probability on every modal-possible event.\\
\bottomrule
\end{longtable}

\section{Audit scope, evidence, and reproducibility}
\subsection{Inputs actually available}
Four mounted ZIP archives were inspected and preserved:
\begin{itemize}
\item \file{Pure_Boolean_Modal_Quantum.zip}
\item \file{Native_Boolean_CM_Rotation_Phase.zip}
\item \file{Native_Boolean_CM_Block_Quantum.zip}
\item \file{Independent_Phase_CM_Tensor_Audit.zip}
\end{itemize}
The user's Windows directory was not accessible. Original archive hashes are recorded in \file{data/input_archives.json}. The original CM manuscript and the revised operator-level manuscript were consulted through File Library excerpts, including the displayed CM/LM definitions and measurement discussion \cite{droncheff2018,theory2026}. Those library PDFs are not silently represented as files read from the Windows directory.

The historical sources are retained under \file{source_snapshots/}. Their original content is evidence to audit, not an authority that makes their interpretations true. All three supplied SHA-256 manifests checked without missing or changed entries: 43, 25, and 19 listed files respectively. The rotation-phase package had no corresponding manifest.

\subsection{What ``exhaustive'' means here}
A \textbf{proved theorem} is justified for its stated mathematical domain. An \textbf{exhaustive finite search} covers the stated finite universe, or uses an exact decision procedure over that universe. A \textbf{regression/spot check} verifies selected instances. A \textbf{literature comparison} is not a novelty proof. An \textbf{open item} has not been established.

All four historical test suites passed again: $22+10+15+12=59$ tests. A new implementation, not importing historical computational modules, recomputed the main inventories, tables, protocols, and counterexamples. A separate 31-test suite checks the new kernel and evidence. The total is \textbf{90 passing named tests}, in addition to the millions of explicit cases in the drivers. A test count is not a substitute for the proofs and domain descriptions below.

The independent kernel uses packed Boolean rows. A second scalar AND/XOR contraction checks the packing; independent rank algorithms are compared on all 65,536 binary $4\times4$ matrices. The contextuality calculation uses bitset implication closure rather than the historical solver, and is cross-checked by brute force in small scenarios. Teleportation is verified both from full tensor contractions and from the derived branch-map formula. The dense-coding calculation is newly recorded as executable evidence rather than merely asserted in prose.

\textbf{Limits of certification.} This audit is not a proof-assistant formalization. It does not exhaust all $2^{64}$ literal bipartite state vectors by enumeration, all $\GL(8,2)$ measurement contexts, or arbitrary ancilla instruments. Universal claims in those larger domains are made only where a proof is given. The earlier signed/complex lift is outside the requested pure-Boolean branch and is not certified here.

\section{Native CM logic and a necessary type ledger}
\subsection{The Boolean contraction}
Let $\B=\{0,1\}$. Every scalar entry in the core calculation is a bit. Define
\begin{equation}
 (A\circ B)_{ij}=\bigoplus_k (A_{ik}\wedge B_{kj}),
 \qquad (A\circ v)_i=\bigoplus_j(A_{ij}\wedge v_j).
 \label{eq:composition}
\end{equation}
The symbol $\circ$ always means this AND/XOR matrix contraction. Pointwise conjunction $A\wedge B$ is a different operation. The notation $\F$ is a mathematical identification of XOR with field addition and AND with field multiplication; it introduces no numerical amplitudes. All ranks below are over this XOR/AND field, \emph{not} the Boolean OR--AND factorization rank sometimes also called ``Boolean rank.''

Integers used as array indices, packed bit masks, counts, dimensions, and external rank diagnostics are not evolving amplitudes. Likewise, the real-valued linear program in Section~\ref{sec:probability} is an external consistency test and is not part of state evolution.

In the original true-first CM convention,
\[
 \ket{x}=\binom{x}{\neg x},\qquad
 \bra{x}C\ket{y}=C_{xy},\qquad x,y\in\{1,0\}.
\]
The contraction selects one cell. There are 16 possible $2\times2$ truth tables. Pointwise AND/XOR/NOT make the same set a Boolean algebra; OR, NAND, NOR, implication, and XNOR are derived pointwise operations. Under XOR and $\circ$ it is instead $M_2(\F)$, with six invertible elements. These are two structures on the same set, not interchangeable compositions.

\subsection{LM extraction does not supply the modal postulate}
A formula-valued LM is a symbolic operator built from the polarity pairs $X,\neg X$ and $Y,\neg Y$. Its evaluated pairing has the identity
\begin{equation}
 \bra{A}[M_{X\Theta Y}]\ket{B}
   =(A\leftrightarrow X)\;\Theta\;(Y\leftrightarrow B).
 \label{eq:lm}
\end{equation}
Valuation commutes with this expression because it preserves AND, XOR, and NOT. The revised manuscript explicitly calls this logical pairing, rather than a physical measurement \cite{theory2026}. The original discussion similarly describes logical truth extraction and does not supply probability densities \cite{droncheff2018}.

The new modal rule ``a nonzero output coefficient is a possible outcome'' is therefore an \emph{additional operational interpretation}. It is not derived from Equation~\eqref{eq:lm}. The same warning applies to selective measurement, discarding, classical outcome messages, and the claim that a coefficient pattern cannot be read nondestructively as an ordinary stored bit string.

\ev{16,384 same-frame truth-fusion cases, 512 signed-axis alignment cases, and 256 general LM pairing valuations independently checked.}

\subsection{Four models, not one silently changing model}
\begin{table}[htbp]
\centering\small
\begin{tabularx}{\textwidth}{@{}L{16mm}L{35mm}Y@{}}
\toprule Model & Objects & Composition and interpretation\\\midrule
\modelC & Native CMs/LMs & Truth tables, formula-valued operators, pointwise Boolean logic, and typed logical pairing.\\
\modelR & Rotation operators $\A$ & Selected XORs of $I_4,P,P^2,P^3$; ordinary Boolean operator composition.\\
\modelS & Shared module $\A^{2^n}$ & One four-coordinate coefficient register for all logical branches; tensor product over $\A$; $\A$-linear operations.\\
\modelL & Literal registers $(\F^8)^{\otimes n}$ & A separate four-position phase factor per party; tensor product over $\F$; restricted or full Boolean-linear gates must be declared.\\
\bottomrule
\end{tabularx}
\caption{For $n=1,2,3$, model \modelS\ has Boolean dimensions $8,16,32$, whereas model \modelL\ has dimensions $8,64,512$.}
\end{table}

In model \modelL, the \emph{restricted} $P$-compatible family consists of $2\times2$ block matrices with blocks in $\A$. The \emph{enlarged} family admits arbitrary local invertible $8\times8$ Boolean matrices and appropriate joint analyzers. Results proved under one family are not automatically results under the other.

\section{Rotation, the derived phase algebra, and interference}
\subsection{Rotation is native; its algebra is a different type from a CM}
Write the four atomic CM positions clockwise as $E_0,E_1,E_2,E_3$, starting at the upper left. Thus $E_0=[\mathrm{AND}]$, $E_1=[\mathrm{UP}]$, $E_2=[\mathrm{NOR}]$, and $E_3=[\mathrm{DOWN}]$. With $v(C)=(c_0,c_1,c_2,c_3)^T$, clockwise rotation is
\begin{equation}
 P=\begin{pmatrix}0&0&0&1\\1&0&0&0\\0&1&0&0\\0&0&1&0\end{pmatrix},
 \quad P^4=I_4,\quad P^2\ne I_4,\quad P^T=P^3.
\end{equation}
Here powers mean repeated $\circ$. Transposition of the underlying CM induces the permutation $K$ interchanging coordinates 1 and 3, and $K\circ P\circ K=P^3$.

Define the operator selected by a CM pattern:
\begin{equation}
 \Lambda(c)=\bigoplus_{j:c_j=1}P^j,
 \qquad \A=\{\Lambda(c):c\in\B^4\}.
 \label{eq:lambda}
\end{equation}
This uses only selection by bits and XOR of actual matrices. Composition is still Equation~\eqref{eq:composition}. Expanding a composition of two such operators gives the cyclic coefficient rule from the earlier work. Thus the previously introduced cyclic product is \emph{derivable}; it is not an extra primitive in the audited implementation.

\begin{theorem}[Exact centralizer]
The matrices commuting with $P$ in $M_4(\F)$ are exactly the 16 operators in $\A$.
\end{theorem}
\begin{proof}
The four vectors $e_0,P e_0,P^2 e_0,P^3 e_0$ form a basis. If $A$ commutes with $P$, its value on $e_0$ determines its values on this whole basis: $A P^j e_0=P^j A e_0$. There are 16 choices for $A e_0$. The 16 matrices in Equation~\eqref{eq:lambda} realize all these choices and commute with $P$.
\end{proof}
\ev{All 65,536 Boolean $4\times4$ matrices scanned; exactly 16 commute. All 256 phase-operator pairs checked with two contraction implementations.}

\subsection{The corrected nilpotent and the limit of the imaginary-number analogy}
Let $N=I_4\oplus P$. Then
\[
 N^2=I_4\oplus P^2,\qquad N^4=0,
 \qquad \rk(N^j)=4-j\quad (0\le j\le4).
\]
The operators $I_4,N,N^2,N^3$ form a basis of $\A$. Consequently $\A$ is isomorphic to $\F[u]/(u^4)$, with $u$ naming $N$ and its product interpreted as $\circ$.

The truth pattern $\delta_{\rm CM}=E_0\oplus E_2=I_2$ satisfies
\[
 \delta_{\rm CM}\wedge\delta_{\rm CM}=\delta_{\rm CM},\qquad
 \delta_{\rm CM}\circ\delta_{\rm CM}=\delta_{\rm CM}.
\]
Its selected rotation operator is instead $\Lambda(\delta_{\rm CM})=N^2$, whose square is zero. No contradiction remains once the types are distinguished. Importantly, $\Lambda$ is \emph{not} a homomorphism for pointwise conjunction or native $2\times2$ composition.

The cycle $I,P,P^2,P^3$ is a phase-clock analogy, not complex scalar multiplication. In characteristic two, additive negation is identical to the original scalar; $P^2$ is not a distinct scalar $-1$. Complement of a coefficient vector is affine, $v\mapsto v\oplus\mathbf1$, and is not complex phase negation.

There is also a correction to a previous minimal-dimension suggestion. No invertible $2\times2$ Boolean matrix has order four, but
\[
 J_3=\begin{pmatrix}1&1&0\\0&1&1\\0&0&1\end{pmatrix}
 \quad\text{satisfies}\quad J_3^4=I_3,\ J_3^2\ne I_3.
\]
Four coordinates are natural for the literal four-position CM permutation; they are not the minimum dimension for an arbitrary order-four Boolean-linear operator.

\subsection{Reversible superposition and phase-sensitive XOR cancellation}
The shear
\[
 C=\begin{pmatrix}1&0\\1&1\end{pmatrix}
 \quad\text{satisfies}\quad C^2=I_2,\qquad C e_0=e_0\oplus e_1.
\]
This is a multibranch vector capable of later cancellation, not a probabilistic mixture. On two four-coordinate branches, use
\[
 C_{\rm sh}=\begin{pmatrix}I_4&0\\I_4&I_4\end{pmatrix},
 \qquad S_k=\begin{pmatrix}I_4&0\\0&P^k\end{pmatrix}.
\]
The circuit $C_{\rm sh}\circ S_k\circ C_{\rm sh}$ maps $(e_0,0)$ to $(e_0,e_0\oplus e_k)$. The second-branch packed patterns are $0,3,5,9$ for $k=0,1,2,3$. Identical paths cancel; distinct rotations leave distinguishable exact patterns. A coarse nonzero detector distinguishes only zero from the other three patterns.

A symmetric splitter is
\[
 H_{\rm rot}=\begin{pmatrix}I_4&N^2\\N^2&I_4\end{pmatrix},
 \qquad H_{\rm rot}^2=I_8,\qquad H_{\rm rot}^T\circ H_{\rm rot}=I_8.
\]
It splits each elementary Boolean coordinate into both logical branches. It does \emph{not} split every nonzero phase pattern: $N^2v=0$ for three nonzero patterns. Twelve of the 15 nonzero phase patterns split. The phase-kickback identity also survives: swapping the branches of $(e_0,e_2)$ equals applying $P^2$ to both branches.

\section{Shared-module structure: what is proved and what is problematic}
\subsection{Units, gates, and Smith classes}
An element of $\A$ is invertible exactly when its selector has odd parity. There are eight units. Reduction modulo the ideal $(N)$ gives the field $\F$. A $2\times2$ block gate is invertible exactly when its reduction is in $\GL(2,2)$, so
\[
 |\GL_2(\A)|=6\cdot8^4=24,576.
\]
Among the 65,536 possible gates, 512 are Boolean-orthogonal after expansion to $8\times8$ matrices. There are 128 monomial invertible gates and 24,448 nonmonomial ones. Nonmonomial here does not mean every basis input splits; it means the gate is not just an invertible block scaling and permutation. Row-unit rescalings and outcome permutations reduce complete invertible two-row bases to 192 unordered modal bases, of which four arise from the orthogonal subgroup.

\begin{theorem}[Two-party rotation-linear normal form]
Every $2\times2$ matrix over $\A$ is equivalent, under independent invertible row and column operations over $\A$, to exactly one
\[
 D_{a,b}=\operatorname{diag}(N^a,N^b),\qquad 0\le a\le b\le4,
\]
where $N^4=0$. There are 15 classes.
\end{theorem}
\begin{proof}
Every nonzero operator is a unit times $N^a$. Choose an entry of minimum valuation, move it to the first diagonal position, and remove its unit factor. Because it divides every other entry, Boolean row/column shears clear its row and column. Reduce the remaining diagonal entry similarly. The ideals generated by entries and the elementary divisor data, equivalently the ranks of successive $N$-filters in this two-generator case, determine the exponents uniquely.
\end{proof}
The expanded Boolean rank profile is
\begin{equation}
 r_j=\max(4-a-j,0)+\max(4-b-j,0),\quad j=0,1,2,3.
 \label{eq:profile}
\end{equation}
All 65,536 matrices were classified independently using these profiles. Table~\ref{tab:classes} records every orbit size.

\subsection{Separability and the scope of gate classification}
A shared two-party state is CM-separable when its $2\times2$ coefficient matrix is an outer product of two $\A$-vectors. This is equivalent to $b=4$ in the normal form. Indeed, a vector over the chain ring can be transformed to $(N^s,0)$ by an invertible change of coordinates; applying this on both outer-product factors leaves at most one nonzero diagonal entry. Conversely, $(N^a,0)$ is explicitly an outer product.

There are 5,266 such vectors including zero, hence 60,270 nonseparable vectors. A zero determinant alone is \emph{not} a separability test over this ring. Local invertibles preserve separability in both directions, since they transform the two factors separately and have local inverses.

The stabilizers of all 15 canonical states were calculated by exact fiber counting, and each satisfies
\[
 |\mathrm{Stab}(D_{a,b})|\,|\mathrm{Orbit}(D_{a,b})|=24,576^2.
\]
The full values are in \file{data/local_stabilizers_audited.csv}. In addition, all 24 two-party basis permutations were tested on all 5,266 separable inputs: 16 create nonseparability for some input; eight preserve it for every input. This is \emph{not} an enumeration of all two-party invertible gates over $\A$.

\subsection{A fundamental preparation-closure failure}
\begin{proposition}[Nonzero factors can have a zero shared tensor]
If every nonzero vector in the shared module is declared an allowed state, independent preparation is not closed under the tensor product over $\A$.
\end{proposition}
\begin{proof}
Take $x=N^3e_0\ne0$ and $y=Ne_0\ne0$ as coefficient vectors on a fixed logical branch. Their shared tensor coefficient is selected by $N^3\circ N=N^4=0$. Thus two nonzero preparations yield the zero joint vector. In the literal tensor product their product is nonzero: a pair of nonzero Boolean coordinates contributes a one.
\end{proof}
In packed phase coordinates this witness is $x=15$, $y=3$. The literal product is nonzero (packed value 13,107), whereas the shared product is zero.

This does \emph{not} invalidate the ring identities or the explicit teleportation contractions. It does invalidate treating the unrestricted collection of all nonzero shared-module vectors as an automatically coherent independent-preparation operational theory. Restricting to primitive/unimodular preparations might repair part of this problem, but would remove many of the previously counted states and requires a separate closure analysis. That repair is not silently assumed here.

\section{Literal tensoring is not an entanglement-preserving change of notation}
\subsection{The Choi-style embedding}
For a shared coefficient matrix $m=(m_{ij})$, form the $8\times8$ Boolean block matrix
\[
 M=\begin{pmatrix}\Lambda(m_{00})&\Lambda(m_{01})\\
 \Lambda(m_{10})&\Lambda(m_{11})\end{pmatrix}.
\]
Its literal bipartite vector is
\begin{equation}
 \ket{M}\!\rangle=\bigoplus_{j,k:M_{jk}=1}\ket{j}_A\otimes\ket{k}_B.
 \label{eq:choi}
\end{equation}
This is an injective linear encoding of 65,536 shared coefficient vectors into a space with $2^{64}$ Boolean vectors. It is not a monoidal equivalence between the two tensor theories.

A literal nonzero bipartite state is a simple tensor exactly when its coefficient matrix has rank one. For the canonical resources,
\[
 \rk M=8-a-b.
\]
Thus the structured Choi family has just nine nonzero literal product states, rather than 5,265; 65,526 of its vectors are literally nonseparable. For example, the shared product state with coefficient matrix $\operatorname{diag}(1,0)$ maps to the literal matrix $\operatorname{diag}(I_4,0)$, of rank four. The encoding has introduced phase correlations. It cannot be used to argue that the two notions of entanglement are identical.

\subsection{The universal resource contains more correlations}
The phase-anchored logical Bell state
\[
 \ket{0,e_0;0,e_0}\oplus\ket{1,e_0;1,e_0}
\]
has literal rank two. The full resource used for universal transfer is
\begin{equation}
 \ket{\Phi_8}=\bigoplus_{j=0}^{7}\ket{j,j},\qquad M=I_8,
 \label{eq:phi8}
\end{equation}
which has rank eight. It includes a correlated four-position phase factor as well as the logical Bell correlation. Increasing the analyzer from four to 64 outcomes is therefore not the only change: the operational resource has changed too.

Equation~\eqref{eq:phi8} is preparable using reversible Boolean gates in the enlarged model. Regard a local label as three binary labels, create a modal superposition on each of Alice's three factors using $C$, and copy the \emph{basis labels} into Bob's initially fixed factors with three controlled-XOR permutations. This is the usual tensor-basis preparation, not cloning an arbitrary unknown modal vector.

\subsection{Local equivalence changes under the larger gate set}
Under full local $\GL(8,2)$, ordinary row/column reduction gives
\[
 U\circ M\circ V=\operatorname{diag}(I_r,0),\qquad r=\rk M.
\]
Hence rank alone classifies literal bipartite pure-state orbits. In particular $(0,2)$ and $(1,1)$ both have rank six and are locally equivalent. Explicit invertible witnesses are saved in \file{data/equal_rank_class_equivalence.json}.

The old distinction between their strong/logical contextuality is therefore a distinction relative to the \emph{restricted 192-basis family}, not an invariant under arbitrary full Boolean-linear local gates. Enlarging the measurement family can expose additional contextuality. No exhaustive full-$\GL(8,2)$ contextuality classification is asserted in this report.

\section{Measurement assumptions, Bell tables, and contextuality}
\subsection{The exact operational contract}
In literal modal quantum theory, a pure state is a \emph{nonzero} vector. A complete measurement is a basis of the dual space, or equivalently an invertible analyzer followed by coordinate readout. An outcome is possible precisely when its evaluated coordinate is one. If only logical outcomes are recorded, the four phase coordinates are coarse-grained: an outcome is possible when its entire output block is nonzero.

Coarse-graining is an \emph{OR of possibilities}, not XOR of outcome amplitudes. These operations must not be confused. A selective outcome leaves the corresponding nonzero conditional vector; disregarding an outcome is naturally represented by the span of conditional vectors. No probabilities are attached to these alternatives. This measurement contract is standard in modal quantum theory \cite{schumacher2011}; its use here is an explicit extension of CM logical extraction.

\begin{proposition}[Possibilistic no-signalling]
For literal tensor states and complete local invertible analyzers, the set of possible outcomes at one party is independent of the other party's choice of local analyzer.
\end{proposition}
\begin{proof}
Fix an Alice effect and contract it into the joint state. The result is a Bob vector (or block of vectors). Bob's invertible basis change cannot transform a nonzero vector to zero or zero to nonzero. Therefore whether at least one Bob outcome is possible is unchanged. The same reasoning works with the parties interchanged and with block coarse-graining.
\end{proof}
This is a support statement. It is weaker than the existence of a no-signalling \emph{probability distribution}; see Section~\ref{sec:probability}.

\subsection{Bell contradiction with only three settings per party}
Use the following ordered dual bases on the logical two-coordinate factor:
\[
 Z=((1,0),(0,1)),\quad X=((1,1),(1,0)),\quad Y=((0,1),(1,1)).
\]
The letters are setting names, not complex Pauli operators. For the support-Bell state, the possible/impossible table is:
\begin{table}[htbp]
\centering
\begin{tabular}{@{}cccc@{}}\toprule
Alice $\backslash$ Bob & $Z$ & $X$ & $Y$\\\midrule
$Z$ & 1001 & 1110 & 0111\\
$X$ & 1110 & 0111 & 1001\\
$Y$ & 0111 & 1001 & 1110\\\bottomrule
\end{tabular}
\caption{Each string lists outcomes $00,01,10,11$; one means possible. Both the shared support-Bell table and the literal correlated construction give this table.}\label{tab:bell}
\end{table}

\begin{theorem}[Explicit modal Bell contradiction]
No deterministic assignment to the six local settings satisfies Table~\ref{tab:bell}.
\end{theorem}
\begin{proof}
The contexts $ZZ,XY,YX$ require
\[
 z=A_Z=B_Z,\quad x=A_X=B_Y,\quad y=A_Y=B_X.
\]
The remaining contexts imply $z\lor x$, $z\lor y$, $\neg(x\wedge z)$, $\neg(y\wedge z)$, $x\lor y$, and $\neg(x\wedge y)$. If $z=0$, then $x=y=1$, contradicting the last condition. If $z=1$, then $x=y=0$, contradicting $x\lor y$.
\end{proof}
\ev{All $2^6=64$ deterministic assignments independently rejected. No expectations or CHSH probabilities are used.}

The contradiction rules out a local relational hidden-variable explanation with simultaneous setting assignments and the stated support semantics. It does not rule out a conventional computer globally calculating these tables, nor derive spatially nonlocal physics from a Boolean truth table.

\subsection{The complete restricted 192-basis result}
A local $\A$-basis consists of two unimodular row effects, reduced by independent multiplication by units and by row permutation. Exactly 192 unordered bases remain. This is a complete family for this \emph{declared block-basis quotient}; it is not every fine-grained literal $8$-outcome measurement.

For each pair of settings, forbidden outcomes produce clauses
\[
 (A_i\ne a)\lor(B_j\ne b).
\]
A global assignment exists exactly when the resulting 2-SAT instance is satisfiable. A possible section extends globally exactly when adding its two outcome literals leaves the instance satisfiable. We use the distinctions of the possibilistic contextuality framework \cite{abramsky2011}:
\begin{itemize}
\item \emph{Strong}: no global assignment exists.
\item \emph{Logical, not strong}: a global assignment exists, but some possible section has no global extension.
\item \emph{Relationally local}: every possible section extends; one global assignment alone would not suffice.
\end{itemize}
The zero state has empty supports and is excluded as a preparation rather than called strongly contextual.

All $192^2=36,864$ setting pairs were computed for every canonical class in both shared and literal models. Exact implication closure decides the constraints over all $2^{384}$ deterministic assignments without pretending to enumerate them one by one. The result is four strong classes, six logical-not-strong classes, four nonzero local classes, and zero. Their structured state counts are 26,214, 34,056, 5,265, and one respectively. The counts and uncovered sections are in Table~\ref{tab:contexts}.

\subsection{Why the shared and literal counts agree}
Define conjugation of a phase operator by ordinary transpose, $\bar A=A^T$. Since $P^T=P^3$, it is an automorphism of $\A$. In a shared contraction the output coefficient is represented by
\[
 \bigoplus_{ij} E_i\circ M_{ij}\circ F_j.
\]
In a literal Choi contraction the corresponding phase block is
\[
 \bigoplus_{ij} E_i\circ M_{ij}\circ F_j^T.
\]
Thus the literal table at $(E,F)$ is the shared table at $(E,\bar F)$, with the row-unit representatives and outcome order adjusted. Conjugation maps the complete 192-basis family to itself.

\begin{proposition}[Empirical-model isomorphism]
The two restricted support models are exactly isomorphic by a permutation of Bob's settings and possible outcome swaps. Equal aggregate contextuality counts are a consequence of this relabeling.
\end{proposition}
\ev{Every context in all 15 classes checked under the explicit relabeling; 8,640 direct literal effect-pair matrix contractions also verify the formula. Brute-force tests on the four-basis family check the 2-SAT coverage algorithm.}

\section{GHZ: exact successes and retained negatives}
The literal correlated three-party state is
\[
 \ket{\mathrm{GHZ}_8}=\bigoplus_{j=0}^{7}\ket{j,j,j}.
\]
Each bipartition flattening has rank eight, so it is nonseparable across all three cuts. This is not a claim about a density-matrix pairwise reduction; the theory here has no such density matrix.

The 27 logical $Z/X/Y$ contexts contain 112 possible sections. All 512 deterministic assignments to the nine local settings fail. A six-context contradiction is
\[
 ZZZ,\quad ZXX,\quad XZY,\quad XXX,\quad YYZ,\quad YYY.
\]
No subset of one through five contexts from this particular 27-context family is contradictory. The exact subset counts tested were 27, 351, 2,925, 17,550, and 80,730: 101,583 subsets in total. Six is therefore minimal \emph{within this family}, not among all possible CM measurements or all possible GHZ constructions.

The four tested orthogonal block bases, by contrast, yield 1,024 global assignments and no uncovered possible sections for this state. Nonseparability is not sufficient to force contextuality in every restricted measurement family.

For a weighted literal variant, apply $N^2$ to Charlie's phase leg in the logical $111$ term, leaving the $000$ term unweighted. Its three flattening ranks are $8,8,6$. The $Z/X/Y$ test is local (16 global assignments); the four orthogonal bases give logical contextuality (291 global assignments and 28 uncovered sections). This is one explicitly specified three-party lift, not a unique Choi lift of the old shared weighted GHZ. The numerical result is verified; selecting a canonical higher-party lifting principle remains a modeling question.

\section{Universal exact teleportation: construction and proof}
\subsection{What ``universal'' requires}
Fix a known shared resource, an input-independent analyzer, and an outcome-dependent correction. Universal exact transfer means that every \emph{nonzero} unknown input vector is recovered exactly on every possible complete outcome. The null input may be tested as a linear identity but is not an additional physical/modal state. No quotient by a global phase or unit is needed for the successful protocols.

We analyze the standard one-copy architecture: Alice holds the unknown $d$-dimensional system and her half of a resource, performs a complete rank-one $d^2$-outcome modal analysis, reports its label, and Bob applies a linear correction. Ancillas and filters will be treated separately rather than hidden in this definition.

\subsection{The conditional branch formula}
Let the resource have coefficient matrix $M$, and reshape analyzer row $m$ into a $d\times d$ matrix $R_m$. With input $\psi$, the coefficient of Bob coordinate $k$ is
\[
 \bigoplus_{i,j}\bigl(\psi_i\wedge (R_m)_{ij}\wedge M_{jk}\bigr).
\]
Consequently the conditional vector is
\begin{equation}
 T_m\circ\psi,\qquad T_m=M^T\circ R_m^T.
 \label{eq:branch}
\end{equation}
This equation is a Boolean contraction identity, not an assumed measurement probability formula. A complete analyzer is invertible when the flattened $R_m$ form a basis of the entire matrix space.

\begin{theorem}[Literal one-copy resource criterion]\label{thm:tele}
In the above literal Boolean-linear architecture, universal exact transfer exists for a $d\times d$ resource whenever an invertible matrix-effect basis exists and $\rk M=d$. Conversely, a resource with rank less than $d$ has no nonzero universally exact success branch. For $d=8$, an explicit invertible effect basis is given below, so rank eight is necessary and sufficient.
\end{theorem}
\begin{proof}
If $C_m\circ T_m=I_d$, then $d=\rk(C_m\circ T_m)\le\rk T_m\le\rk M$. For the converse, choose a complete basis with every $R_m$ invertible. If $M$ is invertible, every $T_m$ is invertible and $C_m=T_m^{-1}$ recovers every input.

Postselecting a nonzero singular branch cannot evade the obstruction. Choose $k\ne0$ in its kernel and $x$ outside its kernel. Both $x$ and $x\oplus k$ give the same nonzero branch output; no fixed correction can recover both distinct inputs. A zero branch is never a success event.
\end{proof}

\subsection{A small, explicit, search-free Bell analyzer}
Use four invertible $2\times2$ effect matrices
\begin{equation}
 L_0=\begin{pmatrix}0&1\\1&1\end{pmatrix},\quad
 L_1=\begin{pmatrix}1&0\\1&1\end{pmatrix},\quad
 L_2=\begin{pmatrix}0&1\\1&0\end{pmatrix},\quad
 L_3=I_2.
 \label{eq:L}
\end{equation}
Their flattened rows form the invertible analyzer
\begin{equation}
 A_4=\begin{pmatrix}0&1&1&1\\1&0&1&1\\0&1&1&0\\1&0&0&1\end{pmatrix}.
\end{equation}
The original modal literature already contains equivalent four-state Boolean Bell constructions and the corresponding teleportation/dense-coding protocols \cite{schumacher2011}; the present audit does not claim their discovery.

For $d=8=2^3$, set
\begin{equation}
 R_{ijk}=L_i\otimes L_j\otimes L_k,\qquad i,j,k\in\{0,1,2,3\}.
 \label{eq:64basis}
\end{equation}
Each matrix is invertible. Since the four $L_i$ form a basis of $M_2(\F)$, their triple tensors form a basis of $M_8(\F)$. The 64-row analyzer with row $\vecop(R_{ijk})^T$ is therefore invertible. This is a tensor construction of three mobit protocols, up to wiring permutations; an unexplained numerical search is unnecessary.

The phase effects $L_j\otimes L_k$ can be synthesized using $P$, its inverse, and
\[
 T=\begin{pmatrix}1&1&0&0\\0&1&0&0\\0&0&1&0\\0&0&0&1\end{pmatrix}.
\]
Exact group generation gives $\langle P,T\rangle=\GL(4,2)$, with 20,160 elements. Explicit words for the 16 canonical phase effects and independent checks of the old 16 synthesis words are included. Crucially, $T\circ P\ne P\circ T$. This is a genuine enlargement beyond the commuting phase algebra. The analyzer is an allowed joint operation on Alice's systems, not a product of disconnected local actions on them.

\subsection{What was executed, not merely inferred}
The fresh code constructs the full $512\times8$ input-to-joint-state preparation and the full $512\times512$ analyzer acting on Alice's pair and Bob's identity factor. It extracts Bob's 64 conditional $8\times8$ maps and checks Equation~\eqref{eq:branch}.

For both the historical literal analyzer and the new tensor-product analyzer, all $256\times64=16,384$ input/outcome identities recover exactly. Of these, \textbf{16,320 per analyzer} concern nonzero inputs and possible outcomes; 64 are null-vector regressions. The old raw received/corrected outputs also match the new independent execution exactly.

For all 65,536 structured resources, all 64 canonical branch ranks were computed explicitly: \textbf{4,194,304 rank comparisons}, with no mismatches. The equality is also proved from invertibility of $R_m$:
\[
 \rk T_m=\rk M.
\]
Direct full-tensor contraction additionally checks all 15 canonical resources and 16 fixed-seed nonsymmetric resources, giving 1,984 branch-formula comparisons. Exactly 24,576 structured resources have rank eight. The rank criterion itself applies to arbitrary $8\times8$ Boolean resources, not only this structured subset; that broader statement is proved, not exhaustively enumerated.

Finally, $C_m\circ T_m=I_8$ implies
\[
 (C_m\circ T_m)\otimes I_r=I_{8r}
\]
for any reference dimension $r$. Thus the identity preserves arbitrary correlations with a reference system. All 64 maps were also explicitly checked with a two-dimensional reference. A full-rank resource is not merely correct on basis states.

\subsection{The older four-outcome shared protocol}
The same logical analyzer acts in the shared coefficient module with four outcomes and recovers all coefficient pairs $(A,B)\in\A^2$, as an exact module identity. All 1,024 input/outcome identities were checked. This includes four null-input regressions. Its operational interpretation retains the preparation-closure limitation identified above.

In the literal model, applying only this logical analyzer leaves 16 separate Alice phase-pair sectors under each reported logical label. All 16 sector maps are nonzero, distinct, and of rank two. A four-valued logical label therefore does not identify a single universally invertible Bob map. The literal construction requires a full 64-outcome phase-and-logical analysis in this architecture.

\section{Dense coding, readout, and resource accounting}
Use $\ket{\Phi_8}$ and the 64 invertible encodings $R_m$ of Equation~\eqref{eq:64basis}. Alice acts only on her subsystem, yielding
\[
 (R_m\otimes I_8)\circ\ket{\Phi_8}=\ket{R_m}\!\rangle.
\]
These 64 joint vectors are linearly independent. Let $Q$ be the matrix with these vectors as columns. After Alice transmits her subsystem, Bob applies $Q^{-1}$ and obtains the unique coordinate vector $e_m$. All 64 states were constructed explicitly, their Boolean rank is 64, and all 64 messages decode exactly.

The dense decoder is $Q^{-1}$; it should not be confused with the teleportation analyzer, whose \emph{rows} are the flattened effects. Their relationship depends on transpose and dual-basis conventions.

Under a single complete modal basis readout, an unassisted $d$-dimensional carrier can perfectly distinguish at most $d$ preparations. If effects distinguish preparations $\psi_j$, one can select functionals $f_i$ with $f_i(\psi_j)=0$ for $i\ne j$ and $f_i(\psi_i)=1$; this forces linear independence. The constructed assisted protocol distinguishes $d^2$ preparations. Thus the $8$ versus $64$ zero-error basis-readout comparison is meaningful \emph{inside the specified modal resource model}.

It is not a Shannon-capacity theorem, a physical bandwidth result, or compression of an eight-bit classical memory into three physical bits. An eight-dimensional Boolean modal vector is stored by eight ordinary coefficients in a conventional simulation. Reading all those coefficients as classical data is a different access model. Labels for 64 outcomes require six classical binary symbols; this is label accounting, not a probability law.

With the restricted $P$-compatible encodings alone, the invertible encodings span only a 16-dimensional subspace of $M_8(\F)$, not a 64-dimensional one. Therefore 64 independent codewords cannot be obtained by that restricted encoding family acting on $\Phi_8$. The full dense-coding result depends on the declared gate enlargement.

\section{Restricted states, quotients, ancillas, and copies}
\subsection{The shared hierarchy remains a restricted theorem}
For $D_{a,b}$ over $\A$, the canonical coefficient map retains the quotient
\begin{equation}
 Q_{a,b}=\A/(N^{4-a})\ \oplus\ \A/(N^{4-b}),
 \qquad |Q_{a,b}|=2^{8-a-b}.
 \label{eq:quotient}
\end{equation}
Here $\A/(N^4)=\A$ and $\A/(N^0)$ is the zero module. Equivalently, this is the quotient by the kernel of the canonical resource action. Cardinality is an external algebraic diagnostic, not physically communicable bits.

The largest common exact family in the one-branch $\A$-linear architecture has 256 vectors for $(0,0)$, 16 for $(0,b)$ with $b>0$, and only the zero vector for $a>0$. The same maxima can be attained by complete four-outcome analyzers for the free-line cases, with reversible corrections.

For the upper bounds, if $a>0$, every entry of the corrected branch lies in $(N)$; such a matrix is nilpotent, so it cannot fix a nonzero vector. If $a=0<b$, the branch has residue rank at most one. Thus $I-T$ has some unit entry, allowing at most one freely chosen $\A$ coordinate in its kernel, hence at most 16 fixed vectors. For achievability take the line $\{(A,0):A\in\A\}$ and the four effect rows
\[
 (1,0,0,0),\ (1,1,0,0),\ (1,0,1,0),\ (1,0,0,1).
\]
They form an invertible analyzer. Outcome-dependent lower shears cancel the extra second component without altering $A$. This also repairs an earlier example that used a projection as a correction when reversibility was intended.

For \emph{fixed-quotient} transfer, require the corrected output to be exactly $D_{a,b}\psi$, not merely an isomorphic outcome-dependent quotient. The maximum number of independent successful analyzer rows is four when $a=b$, two when $a<b<4$, and three when $a<4,b=4$. The zero case is a trivial zero output, not transmission. These maxima allow arbitrary $\A$-linear corrections; they should not be silently promoted to an all-reversible correction theorem.

\ev{For each of 15 classes, all 65,536 corrections were tested for fixed vectors and all 65,536 analyzer rows were tested for the fixed-quotient image-inclusion condition: 983,040 cases of each type. The rank of their residue rows gives the maximum number that can be part of an invertible analyzer. A separate 256-case check verifies the free-line construction with reversible corrections.}

\subsection{The literal restricted hierarchy is different}
\begin{theorem}[Maximum literal exact subspace]
For a known literal resource of rank $r$, under unrestricted Boolean-linear analyzers and reversible receiver corrections, the largest common exactly transferable linear subspace has dimension $r$. A complete 64-outcome protocol attaining it exists for each $1\le r\le8$.
\end{theorem}
\begin{proof}
A branch that is the identity on a subspace $W$ is injective there, so $\dim W\le\rk T\le r$. Reduce the resource locally to $D=\operatorname{diag}(I_r,0)$. Choose 64 linearly independent effects whose upper-left $r\times r$ blocks are invertible. Such a basis exists: invertible $r\times r$ matrices span $M_r(\F)$ (the audit gives explicit identities, transvections, and permutation variants); embed a basis of them, then use $I_r$ with one added entry at a time outside that corner. These span all 64 matrix directions while retaining an invertible corner. On $W=\spn(e_0,\ldots,e_{r-1})$, each branch is invertible onto that same subspace. Extend its inverse to an invertible operation on all eight coordinates.
\end{proof}
All input/outcome cases in these constructed subspaces were tested: 32,640 identities for $r=1,\ldots,8$, including zero regressions. A rank-six resource such as $(0,2)$ or $(1,1)$ therefore supports a 64-vector literal exact subspace (63 actual nonzero modal states), not the old shared bounds of 16 or one.

A branch rank of six also gives a six-dimensional quotient of the full input space. Equal ranks do \emph{not} alone prove that every outcome has the same kernel or that one fixed quotient encoding is recovered. Statements about a fixed quotient need the explicit corrected-map equation used in the previous subsection.

\subsection{Single-copy postselection and product ancillas}
For a singular resource, Theorem~\ref{thm:tele} rules out even one nonzero universally exact branch. Local product ancillas and local linear processing do not increase bipartite coefficient-matrix rank. Thus they do not make a single rank-six resource transmit an arbitrary eight-dimensional input. Heralding a restricted or quotient transfer is a different task from heralding exact recovery of every unknown input.

No success probability is assigned. ``Heralded'' means that a designated outcome is possible and its conditional output has the stated property; it does not mean the event occurs with a particular frequency.

\subsection{Finite-copy activation: the old no-go cannot be exported}
\begin{theorem}[Literal heralded activation by rank]
With literal tensor products, arbitrary local invertible Boolean-linear maps, and coordinate filters, $k$ copies of a rank-$r$ resource can yield $\Phi_8$ on a heralded branch exactly when $r^k\ge8$.
\end{theorem}
\begin{proof}
Tensor rank for coefficient matrices multiplies, so the $k$-copy resource has rank $r^k$. Local filters cannot increase rank, giving necessity. For sufficiency reduce the resource by local invertibles to $\operatorname{diag}(I_{r^k},0)$, then project each local system onto the first eight coordinates. The resulting coefficient matrix is $I_8$ and is nonzero, so the designated branch is possible.
\end{proof}
For $(0,2)$, two copies have rank $6^2=36$ on a $64\times64$ coefficient matrix. The supplied exact filter matrices satisfy
\[
 F_A\circ(M_{02}\otimes M_{02})\circ F_B^T=I_8.
\]
The witness is stored in \file{data/literal_activation_witness.json}. Ranks 3 through 7 need two copies, rank two needs three, and rank one never suffices. A full-rank resource already works with one copy. The zero resource is not a preparation.

By contrast, in the shared $\A$-linear architecture, reducing a singular resource modulo $(N)$ gives rank at most one over $\F$. Its tensor powers still have residue rank at most one, and product ancillas do not increase it. It cannot yield an identity on a free two-dimensional $\A$ input. That no-activation theorem is valid \emph{in that restricted tensor category}. The two conclusions are compatible because the tensor products and admissible filters are different.

The literal activation theorem proves \emph{heralded extraction}; it does not establish deterministic extraction or a numerical success probability. Minimum copy counts in the tables refer to precisely this task.

\section{Orthogonality and the proposed symplectic repair}
The successful analyzers are invertible but generally not Boolean-orthogonal. For the canonical $d=8$ protocol, the 64-outcome analyzer fails $A^T\circ A=I$, and only eight of its 64 corrections satisfy transpose-orthogonality. At the logical $d=2$ level, two of four corrections are orthogonal.

An exhaustive restricted search over all 65,536 possible four-block outcome rows found 16,384 satisfying the necessary orthogonal-row condition $E\circ E^T=I_4$. Of those, 8,192 give invertible receiver branches, and \emph{none} gives an orthogonally correctable branch for the support-Bell resource. Satisfying the row condition is necessary, not a proof of completion to an orthogonal analyzer.

The natural shared block Bell analyzer obtained by inverting the $H_{\rm rot}$--CNOT preparation circuit was also checked explicitly. Its four branch ranks are $6,6,6,6$ for the support-Bell resource and $4,4,4,4$ for $(0,2)$. The general reversible analyzer instead preserves ranks eight and six respectively. Thus the natural orthogonal splitter is valid, but its associated analyzer discards additional information and is not an optimal partial-transfer construction.


\begin{theorem}[Even-dimensional orthogonal-correction obstruction]
For even $d$, no complete rank-one $d^2$-outcome Bell analyzer can have all branches exactly corrected by transpose-orthogonal matrices, for any invertible resource $M$.
\end{theorem}
\begin{proof}
If $O^T\circ O=I_d$, every column has odd parity. Since $d$ is even, the parity of all entries of $O$ is even. Therefore all flattened orthogonal matrices lie in a proper hyperplane of $M_d(\F)$ and cannot span its $d^2$ dimensions.

If $C_m$ is orthogonal and $C_m\circ T_m=I_d$, then $T_m=C_m^{-1}$ is orthogonal. Equation~\eqref{eq:branch} implies $R_m^T=(M^T)^{-1}\circ T_m$. An invertible linear image of a proper hyperplane is still proper, so the $R_m$ cannot form a complete basis. This contradicts analyzer invertibility.
\end{proof}
The theorem does not cover every conceivable ancilla-assisted instrument or alternative bilinear pairing. It is not a general impossibility theorem for all characteristic-two analogues of unitarity.

However, simply replacing the Euclidean transpose form by a symplectic form cannot preserve the \emph{entire} currently enlarged $P,T$ gate set. Every Boolean $4\times4$ bilinear matrix $J$ satisfying both $P^T\circ J\circ P=J$ and $T^T\circ J\circ T=J$ was enumerated; only $J=0$ survives. Therefore no nondegenerate alternating or symmetric form makes both generators structure-preserving. A form-preserving subtheory must change the gate set, representation, or operational construction. Binary symplectic descriptions of ordinary stabilizer quantum mechanics act on Pauli labels and should not be identified with these Boolean amplitude vectors \cite{gottesman1997}.

\section{No-cloning, exact readout, and what Boolean logic can still do}
An arbitrary Boolean truth function $f$ can be put into a reversible basis-label oracle,
\[
 \ket{x,b}\longmapsto\ket{x,b\oplus f(x)}.
\]
Even when $f$ is nonlinear in the label bits, this is a fixed permutation and therefore acts linearly on the modal coefficient vector. This is different from applying a nonlinear operation directly to the unknown coefficients. Allowing arbitrary Boolean connectives in an oracle does not by itself violate the linear dynamics contract; allowing arbitrary coefficient-dependent gates does.

\begin{theorem}[Linear modal no-cloning]
For dimension at least two, no linear map can send every nonzero $v$ to $v\otimes v$, even when a fixed blank system is included.
\end{theorem}
\begin{proof}
If two independent basis states $e_0,e_1$ are cloned, linearity sends $e_0\oplus e_1$ to $e_0\otimes e_0\oplus e_1\otimes e_1$. Its tensor square also has the distinct cross terms $e_0\otimes e_1$ and $e_1\otimes e_0$. They do not cancel. Including a machine with input-dependent final states does not create these missing tensor-coordinate sectors.
\end{proof}
The basis-copying linear map in dimension eight fails on all 247 nonzero vectors with support size at least two. It works on the eight basis vectors and trivially on zero. This is a regression illustrating the theorem, not its only justification.

It is nevertheless possible to copy a \emph{stored classical coefficient string}: $(v,0)\mapsto(v,v)$ is a reversible Boolean shear on a direct sum. It is also possible to compute the coefficients $v_i\wedge v_j$ using ordinary nonlinear Boolean logic and thus build the array for $v\otimes v$. The latter map is not XOR-linear in $v$ and is outside the declared modal evolution class.

Consequently, no-cloning and the modal communication interpretation require both a gate restriction and an access restriction. An implementation that exposes all CM coefficients as ordinary readable memory is a classical simulation of the model. It does not inherit physical no-cloning, secrecy, or communication advantage merely because its matrices use the same symbols.

\section{A probability obstruction, not just an unspecified Born rule}\label{sec:probability}
The Boolean model does not assign probabilities. One can nevertheless ask, externally, whether Table~\ref{tab:bell} can be represented by ordinary nonnegative probabilities satisfying normalization, no-signalling, and
\[
 p(a,b\mid X,Y)>0\quad\Longleftrightarrow\quad
 (a,b\mid X,Y)\text{ is modal-possible}.
\]
The answer is no. This obstruction is known in modal quantum theory \cite{schumacher2010}; it was independently recomputed for the exact convention used here.

\begin{theorem}[No support-faithful no-signalling probability extension]
Every normalized no-signalling nonnegative probability assignment respecting the forbidden events of Table~\ref{tab:bell} assigns probability zero to six events that the modal table calls possible.
\end{theorem}
\begin{proof}
Solve the linear equalities for normalization, no-signalling, and the forbidden events. Write $x=p(11\mid YZ)$, $y=p(01\mid YY)$, and $z=p(10\mid YY)$. The equalities imply
\begin{align*}
 p(00\mid ZX)&=-x, & p(11\mid YZ)&=x,\\
 p(11\mid XX)&=y+z-1, & p(00\mid YY)&=1-y-z,\\
 p(11\mid ZY)&=x+y-z, & p(00\mid XZ)&=-x-y+z.
\end{align*}
Nonnegativity of each displayed pair forces $x=0$, $y+z=1$, and $y=z$. Hence $y=z=1/2$ and all six displayed events have probability zero. Each is marked possible in the Boolean table.
\end{proof}
These equations use ordinary real arithmetic \emph{only as an external diagnostic}. No negative or fractional value has been inserted into CM state evolution. The resulting weak resolution assigns $1/2$ to the other allowed events, but deletes six modal possibilities, so it is not a faithful probability realization of the original supports.

\ev{Linear programming separately maximized every allowed event and the common minimum allowed probability; an exact symbolic affine solution supplies the proof above. Raw results and exact constraints are included.}

Thus a physical interpretation cannot simply append an ordinary no-signalling probability rule while retaining all these predictions unchanged. It would need to restrict states or measurements, alter the modal interpretation, or otherwise change the operational model. This is stronger than saying that a Born rule has not yet been discovered.

\section{Consolidated class and evidence tables}
\subsection{All 15 structured classes}
\begin{longtable}{@{}ccrcccc@{}}
\caption{Complete structured family. The shared separability column refers to model \modelS; literal nonzero product states occur only at rank one. The zero vector is bookkeeping, not an allowed modal preparation.}\label{tab:classes}\\
\toprule $a$ & $b$ & Orbit size & Literal rank & Shared sep. & 192-basis status & Universal\\\midrule\endfirsthead
\toprule $a$ & $b$ & Orbit size & Literal rank & Shared sep. & 192-basis status & Universal\\\midrule\endhead
0 & 0 & 24,576 & 8 & No & Strong & Yes\\
0 & 1 & 18,432 & 7 & No & Logical & No\\
0 & 2 & 9,216 & 6 & No & Logical & No\\
0 & 3 & 4,608 & 5 & No & Logical & No\\
0 & 4 & 4,608 & 4 & Yes & Local & No\\
1 & 1 & 1,536 & 6 & No & Strong & No\\
1 & 2 & 1,152 & 5 & No & Logical & No\\
1 & 3 & 576 & 4 & No & Logical & No\\
1 & 4 & 576 & 3 & Yes & Local & No\\
2 & 2 & 96 & 4 & No & Strong & No\\
2 & 3 & 72 & 3 & No & Logical & No\\
2 & 4 & 72 & 2 & Yes & Local & No\\
3 & 3 & 6 & 2 & No & Strong & No\\
3 & 4 & 9 & 1 & Yes & Local & No\\
4 & 4 & 1 & 0 & Yes & Excluded & No\\
\bottomrule\end{longtable}
\subsection{Exact contextuality coverage counts}
\begin{longtable}{@{}ccrrrc@{}}
\caption{Per-class counts in the complete 192-basis family, identical under the shared/literal setting relabeling. ``Extendable'' counts allowed joint sections having some compatible global assignment.}\label{tab:contexts}\\
\toprule $a$ & $b$ & Possible & Extendable & Uncovered & Status\\\midrule\endfirsthead
\toprule $a$ & $b$ & Possible & Extendable & Uncovered & Status\\\midrule\endhead
0 & 0 & 141,312 & 0 & 141,312 & Strong\\
0 & 1 & 139,264 & 65,536 & 73,728 & Logical\\
0 & 2 & 137,216 & 100,352 & 36,864 & Logical\\
0 & 3 & 135,168 & 119,808 & 15,360 & Logical\\
0 & 4 & 131,072 & 131,072 & 0 & Local\\
1 & 1 & 135,168 & 0 & 135,168 & Strong\\
1 & 2 & 131,072 & 65,536 & 65,536 & Logical\\
1 & 3 & 126,976 & 98,304 & 28,672 & Logical\\
1 & 4 & 118,784 & 118,784 & 0 & Local\\
2 & 2 & 122,880 & 0 & 122,880 & Strong\\
2 & 3 & 114,688 & 65,536 & 49,152 & Logical\\
2 & 4 & 98,304 & 98,304 & 0 & Local\\
3 & 3 & 98,304 & 0 & 98,304 & Strong\\
3 & 4 & 65,536 & 65,536 & 0 & Local\\
4 & 4 & 0 & 0 & 0 & Excluded\\
\bottomrule\end{longtable}
\subsection{Do not merge the two resource hierarchies}
\begin{longtable}{@{}ccrrrr@{}}
\caption{Largest exact families count vectors including zero. Subtract one for nonzero modal preparations. The last column is the minimum literal-copy count for \emph{heralded} extraction of $I_8$ under unrestricted local Boolean-linear processing.}\label{tab:hierarchy}\\
\toprule $a$ & $b$ & Rank & Shared exact & Literal exact & Literal copies\\\midrule\endfirsthead
\toprule $a$ & $b$ & Rank & Shared exact & Literal exact & Literal copies\\\midrule\endhead
0 & 0 & 8 & 256 & 256 & 1\\
0 & 1 & 7 & 16 & 128 & 2\\
0 & 2 & 6 & 16 & 64 & 2\\
0 & 3 & 5 & 16 & 32 & 2\\
0 & 4 & 4 & 16 & 16 & 2\\
1 & 1 & 6 & 1 & 64 & 2\\
1 & 2 & 5 & 1 & 32 & 2\\
1 & 3 & 4 & 1 & 16 & 2\\
1 & 4 & 3 & 1 & 8 & 2\\
2 & 2 & 4 & 1 & 16 & 2\\
2 & 3 & 3 & 1 & 8 & 2\\
2 & 4 & 2 & 1 & 4 & 3\\
3 & 3 & 2 & 1 & 4 & 3\\
3 & 4 & 1 & 1 & 2 & never\\
4 & 4 & 0 & 1 & 1 & never\\
\bottomrule\end{longtable}
\subsection{Evidence ledger}
\small
\begin{longtable}{@{}L{29mm}L{53mm}L{21mm}L{45mm}@{}}
\caption{Executed universes and proof scopes. All counts are diagnostics, not amplitude values.}\\
\toprule Item & Universe or evidence & Result & Boundary\\\midrule\endfirsthead
\toprule Item & Universe or evidence & Result & Boundary\\\midrule\endhead
CM/LM identities & 16,384 truth-fusion, 512 alignment, 256 pairing cases & Confirmed & Native Boolean logic, not modal measurement.\\
Packed Boolean kernel & All 256 pairs of 2x2 matrices; all 65,536 4x4 ranks and inverses; rectangular checks & Confirmed & Independent scalar and packed implementations.\\
Rotation centralizer & All 65,536 binary 4x4 matrices & 16 operators & Exactly the XOR span of I, P, P squared, P cubed.\\
Rotation powers & Exact identities and ranks 4,3,2,1,0 & Confirmed & Different types for the CM identity and its selected operator.\\
Order-four minimum & All six GL2 maps; explicit 3x3 witness & Correction & Four is not the minimum arbitrary linear dimension.\\
Gate inventory & All 65,536 phase-block gates & 24,576 / 512 & Invertible / Boolean-orthogonal; 128 monomial.\\
Two-party separability & All 65,536 coefficient states and all 65,536 factor pairs & Confirmed & 5,266 shared products including zero, versus 10 literal structured vectors of rank at most one.\\
Local stabilizers & All 15 classes; exact relation-fiber counts & Confirmed & Orbit-stabilizer equality in GL2(A) squared.\\
Basis-permutation gates & 24 permutations on all 5,266 shared product vectors & 16 entanglers & Not all possible two-party gates.\\
Restricted contextuality & 15 classes x 36,864 contexts x two models & Confirmed & Exact implication-closure coverage; not brute enumeration of 2\textsuperscript{384} assignments.\\
Shared/literal relation & Every restricted context; 8,640 direct effect-pair checks & Relabeling & Bob coefficient conjugation explains all table changes.\\
Bell & Nine contexts, all 64 assignments & Contradiction & Modal local-assignment assumptions explicitly stated.\\
GHZ & 27 contexts, all 512 assignments; 101,583 subsets below six & Minimum six & Only this stated context family.\\
Weighted GHZ & One specified literal lift; Z/X/Y and four orthogonal bases & Mixed & Local in the first family; logically contextual in the second.\\
P and shear generation & Exact closure of P, inverse P, and T & 20,160 & All GL4; T leaves the commuting phase subalgebra.\\
Shared teleportation & 256 inputs x four outcomes & 1,024 identities & Includes four null regressions; module preparation caveat remains.\\
Literal teleportation & Two analyzers, each 256 inputs x 64 outcomes & 32,768 identities & 16,320 meaningful nonzero cases per analyzer.\\
Direct branch formula & 15 canonical plus 16 fixed-seed resources, all 64 branches & 1,984 checks & Uses explicit full tensor matrices, including nonsymmetric resources.\\
All resource ranks & 65,536 resources x 64 branches & 4,194,304 checks & No branch-rank mismatch.\\
Reference preservation & 64 corrected identity maps tensor I2 & Confirmed & Proof extends to every reference dimension.\\
Dense coding & 64 encoded joint states and full inverse decoder & 64/64 & Basis rank 64; not a physical bandwidth claim.\\
Old logical-only analyzer & Four labels, each with 16 phase sectors & Fails & Sixteen distinct rank-two Bob maps per logical label.\\
Shared restricted transfer & 983,040 correction fixed-space cases & Confirmed & 256, 16, or one vector depending on class.\\
Shared fixed quotient & 983,040 analyzer-row solvability tests & 4 / 2 / 3 labels & Scope permits arbitrary A-linear corrections; zero quotient is trivial.\\
Literal restricted transfer & Rank 1 through 8, complete analyzers & 32,640 checks & Attains all $2^r$ vectors in an r-dimensional subspace.\\
Literal finite-copy activation & All 15 canonical classes; explicit filters & Positive & Two copies of rank six can heraldedly yield rank eight.\\
Shared tensor closure & Explicit two nonzero factors & Fails & N cubed times N is zero under the shared tensor.\\
Shared orthogonal rows & All 65,536 four-block rows & Zero correctable & Among rows meeting the necessary orthogonal condition.\\
Orthogonal Bell basis & Even-dimensional hyperplane proof & No-go & Complete rank-one analyzer with all corrections orthogonal.\\
Common bilinear form & All 65,536 four-dimensional forms & Only zero & No nondegenerate form preserves both P and T.\\
No-cloning & Proof and all 256 basis-cloner inputs & Linear no-go & 247 multi-support failures; coefficient copying remains possible.\\
Probability extension & External real LP and exact affine solution & Fails faithfully & Six modal-possible events forced to zero by no-signalling.\\
Historical regression & 22 + 10 + 15 + 12 named tests & 59 pass & Passing old tests does not certify their interpretation.\\
New test suite & 31 named independent tests & 31 pass & In addition to all driver enumerations.\\
\bottomrule\end{longtable}\normalsize


\section{Literature, novelty, and useful research directions}
Modal quantum theory already supplies finite-field state spaces, general invertible evolution, support-based measurements, entangled states, no-cloning, Bell phenomena, teleportation, and dense coding \cite{schumacher2011,james2011}. The elementary four-state protocol underlying the present $d=8$ construction is not a new quantum-information primitive. The independent tensor construction makes that relationship particularly explicit.

The global-section distinction between logical and strong contextuality is established prior work \cite{abramsky2011}. Generalized/categorical treatments also separate the algebra of protocols from their complex-Hilbert-space realization \cite{abramskycoecke2004}. Werner's classification concerns tight complex quantum teleportation and dense coding with the appropriate unitary resources; the Boolean full-rank criterion must not be imported as a deterministic physical-quantum criterion for arbitrary partially entangled states \cite{werner2001}.

``Quantum mechanics over sets'' uses a related but different formalism and should not be used to supply a probability rule to this model without additional changes \cite{ellerman2013}. Matrix-logic square-root-of-NOT constructions that introduce complex matrices likewise do not make the native Boolean $P$ a complex scalar \cite{mizraji2021}. Finite rings and finite-chain-ring coding constructions already have a literature \cite{liu2017}; no novelty claim is based merely on using a ring with nilpotents.

The best-supported distinct research object is narrower: a typed CM rotation-operator subalgebra, its exact restricted gate and measurement families, the finite classification relative to those families, and the comparison with the enlarged literal modal theory. The fresh corrections themselves are useful: they identify which properties depend on rotation-linearity, a tensor choice, an access assumption, or an extra correlation resource.

Practical uses presently supported by evidence are exact finite-model experimentation, verification examples, symbolic Boolean operator analysis, and educational/foundational separation of assumptions. Coding theory, error correction, algorithmic advantage, physical implementations, and resource-efficient simulation remain prospective applications, not demonstrated outcomes of these tests. A bounded literature search is not a proof of novelty.

\section{What may now be used as a foundation, and what remains open}
\subsection{A safe construction contract}
For future literal-model work, specify a nonzero Boolean state vector, its actual subsystem tensor factors, the allowed fixed Boolean-linear gates, a basis or coarse-grained modal readout, the selective outcome rule, and the classical message/resource accounting. Record explicitly when leaving the rotation-compatible subgroup. Under that contract, the full-rank teleportation theorem, literal nonseparability criterion, and dense-coding construction are established building blocks.

For rotation-linear work, retain $\A$ as a derived operator algebra, not the native truth-pattern composition. Treat its Smith and quotient results as module algebra, and do not declare all nonzero coefficient states independently preparable until the zero-tensor problem has been addressed. Exact coefficient recovery and recovery modulo units or annihilator ideals are different tasks.

\subsection{Open or deliberately unclaimed items}
The present audit does not establish a canonical repair of shared-module preparation closure; a complete contextuality classification under all literal $\GL(8,2)$ fine-grained measurements; a canonical higher-party Choi lift; a deterministic/probabilistic optimality theory for activation; a replacement symplectic protocol satisfying an alternative operational structure; monogamy inequalities; a general Clifford/stabilizer subtheory with physical Pauli interpretation; a developed CM error-correcting code; a new efficient quantum-computing algorithm; or hardware/experimental predictions.

Multiplication of a complete shared state by a unit preserves zero/nonzero support under the matching $\A$-linear readout, but it does not preserve an exact CM pattern. No global-unit quotient is imposed on the successful exact protocols. Under enlarged literal readout, such a transformation is generally a detectable operator, not automatically a harmless global phase.

\textbf{Final assessment.} The positive Boolean protocol identities are real. The earlier tendency to combine them, their counts, and their limitations into one unrestricted ``CM quantum theory'' was not justified. The corrected baseline is a family of explicitly related algebras and modal models. That is a stronger basis for future work than either discarding the verified mathematics or treating every earlier interpretation as confirmed.

\appendix
\section{Reproducing and inspecting the audit}
The archive contains this LaTeX source and compiled PDF under \file{paper/}, the new kernel/drivers under \file{code/}, raw CSV/JSON under \file{data/}, a 31-test independent suite under \file{tests/}, historical source snapshots, and execution logs. It includes a machine-readable claim ledger and a SHA-256 manifest. No font files or credentials are distributed.

From the extracted root, with the listed Python packages installed, run:
\begin{verbatim}
python run_all.py --historical
\end{verbatim}
The historical flag reruns the four preserved old suites as well as the fresh calculations. Without it, the command reruns the independent calculations, evidence validation, and new tests. The report can be compiled without Python:
\begin{verbatim}
cd paper
latexmk -pdf -interaction=nonstopmode CM_Correctness_Audit.tex
\end{verbatim}
Alternatively, run \file{pdflatex} at least three times from a clean directory, until the table of contents and references stabilize. All generated table source is included, so compilation does not require importing a data file through LaTeX. The empirical timing logs are execution receipts, not performance benchmarks.

\subsection{Raw matrix convention}
Matrices in JSON are lists of packed rows, with bit zero representing column zero. A vector bit zero represents coordinate zero. Flattening is row-major: matrix entry $(i,j)$ occupies bit $i n+j$ for an $n$-column matrix. This convention differs from visually printing an ordinary binary integer most-significant-bit first; use the declared bit positions rather than guessing column order from a printed integer.

\subsection{Why the evidence is reproducible but not a universal certification}
The code contains executable assertions, independent kernel cross-checks, exhaustive inventories, exact implication solving, explicit tensor contractions, and mathematical proofs in this document. The old implementations were not imported to make the new numerical answers. Historical data were used for regression comparisons and to verify supplied synthesis words. Nevertheless, these are implementations and human-readable derivations, not machine-checked proofs of all possible claims. The complete ledger distinguishes the verified universes from the open domains.

\begin{thebibliography}{99}
\bibitem{droncheff2018} B. Droncheff. \emph{Correspondence Matrices; Algorithms for Propositional Logic}. Research manuscript, May 2018, version 1.0.1. DOI: \href{https://doi.org/10.13140/RG.2.2.28036.37764}{10.13140/RG.2.2.28036.37764}. Original CM, rotation, XOR, and LM definitions; Sections 3--4 consulted through the supplied File Library material.
\bibitem{theory2026} B. Theory. \emph{Operator-Level Boolean Computation with Correspondence Matrices}. Unpublished manuscript supplied in File Library, uploaded 17 September 2026. Sections 2 and 8 and Appendices A--B: Boolean contraction, formula-valued LMs, logical pairing, and evidence boundaries.
\bibitem{schumacher2011} B. Schumacher and M. D. Westmoreland. \emph{Modal quantum theory}. Foundations of Physics \textbf{42}, 918--925 (2012); arXiv version 2, 2011. \href{https://arxiv.org/abs/1010.2929}{arXiv:1010.2929}. Finite-field modal states, measurements, Bell/no-cloning, and dense-coding/teleportation constructions.
\bibitem{schumacher2010} B. Schumacher and M. D. Westmoreland. \emph{Non-contextuality and free will in modal quantum theory}. 2010. \href{https://arxiv.org/abs/1010.5452}{arXiv:1010.5452}. Modal constraints and the obstruction to a support-faithful no-signalling probability assignment.
\bibitem{james2011} R. P. James, G. Ortiz, and A. Sabry. \emph{Quantum Computing over Finite Fields}. 2011. \href{https://arxiv.org/abs/1101.3764}{arXiv:1101.3764}. Finite-field interference and information-processing analogues.
\bibitem{abramsky2011} S. Abramsky and A. Brandenburger. \emph{The Sheaf-Theoretic Structure of Non-Locality and Contextuality}. New Journal of Physics \textbf{13}, 113036 (2011). \href{https://arxiv.org/abs/1102.0264}{arXiv:1102.0264}. Global sections and possibilistic contextuality.
\bibitem{werner2001} R. F. Werner. \emph{All Teleportation and Dense Coding Schemes}. Journal of Physics A \textbf{34}, 7081--7094 (2001). \href{https://arxiv.org/abs/quant-ph/0003070}{arXiv:quant-ph/0003070}. Tight complex-quantum protocol classification.
\bibitem{abramskycoecke2004} S. Abramsky and B. Coecke. \emph{A categorical semantics of quantum protocols}. 2004. \href{https://arxiv.org/abs/quant-ph/0402130}{arXiv:quant-ph/0402130}.
\bibitem{gottesman1997} D. Gottesman. \emph{Stabilizer Codes and Quantum Error Correction}. PhD thesis, 1997. \href{https://arxiv.org/abs/quant-ph/9705052}{arXiv:quant-ph/9705052}. Stabilizer and binary symplectic machinery, not a Boolean-amplitude replacement for complex quantum mechanics.
\bibitem{ellerman2013} D. Ellerman. \emph{Quantum mechanics over sets}. 2013. \href{https://arxiv.org/abs/1310.8221}{arXiv:1310.8221}. A related but different set-based formalism.
\bibitem{mizraji2021} E. Mizraji. \emph{Generalized ``Square roots of Not'' matrices, their application to the unveiling of hidden logical operators and to the definition of fully matrix circular Euler functions} 2021; revised 2024. \href{https://arxiv.org/abs/2107.06067}{arXiv:2107.06067}. Matrix-logic comparison; no claim of a native Boolean complex unit.
\bibitem{liu2017} X. Liu and H. Liu. \emph{Quantum Codes from Linear Codes over Finite Chain Rings}. 2017. \href{https://arxiv.org/abs/1704.06375}{arXiv:1704.06375}. Prior finite-chain-ring coding constructions, not a derivation of the present CM resource interpretation.
\end{thebibliography}
\end{document}
